\documentclass{amsart}
% For dealing with references we use the comment environment
\usepackage{verbatim}
\newenvironment{reference}{\comment}{\endcomment}
%\newenvironment{reference}{}{}
\newenvironment{slogan}{\comment}{\endcomment}
\newenvironment{history}{\comment}{\endcomment}

% For commutative diagrams we use Xy-pic
\usepackage[all]{xy}

% We use 2cell for 2-commutative diagrams.
\xyoption{2cell}
\UseAllTwocells

% We use multicol for the list of chapters between chapters
\usepackage{multicol}

% This is generally recommended for better output
\usepackage{lmodern}
\usepackage[T1]{fontenc}

% For cross-file-references

% Package for hypertext links:
\usepackage{hyperref}

% For any local file, say "hello.tex" you want to link to please

% Theorem environments.
%
\theoremstyle{plain}
\newtheorem{theorem}[subsection]{Theorem}
\newtheorem{proposition}[subsection]{Proposition}
\newtheorem{lemma}[subsection]{Lemma}

\theoremstyle{definition}
\newtheorem{definition}[subsection]{Definition}
\newtheorem{example}[subsection]{Example}
\newtheorem{exercise}[subsection]{Exercise}
\newtheorem{situation}[subsection]{Situation}

\theoremstyle{remark}
\newtheorem{remark}[subsection]{Remark}
\newtheorem{remarks}[subsection]{Remarks}

\numberwithin{equation}{subsection}

% Macros
%
\def\lim{\mathop{\mathrm{lim}}\nolimits}
\def\colim{\mathop{\mathrm{colim}}\nolimits}
\def\Spec{\mathop{\mathrm{Spec}}}
\def\Hom{\mathop{\mathrm{Hom}}\nolimits}
\def\Ext{\mathop{\mathrm{Ext}}\nolimits}
\def\SheafHom{\mathop{\mathcal{H}\!\mathit{om}}\nolimits}
\def\SheafExt{\mathop{\mathcal{E}\!\mathit{xt}}\nolimits}
\def\Sch{\mathit{Sch}}
\def\Mor{\mathop{\mathrm{Mor}}\nolimits}
\def\Ob{\mathop{\mathrm{Ob}}\nolimits}
\def\Sh{\mathop{\mathit{Sh}}\nolimits}
\def\NL{\mathop{N\!L}\nolimits}
\def\CH{\mathop{\mathrm{CH}}\nolimits}
\def\proetale{{pro\text{-}\acute{e}tale}}
\def\etale{{\acute{e}tale}}
\def\QCoh{\mathit{QCoh}}
\def\Ker{\mathop{\mathrm{Ker}}}
\def\Im{\mathop{\mathrm{Im}}}
\def\Coker{\mathop{\mathrm{Coker}}}
\def\Coim{\mathop{\mathrm{Coim}}}

% Boxtimes
%
\DeclareMathSymbol{\boxtimes}{\mathbin}{AMSa}{"02}

%
% Macros for moduli stacks/spaces
%
\def\QCohstack{\mathcal{QC}\!\mathit{oh}}
\def\Cohstack{\mathcal{C}\!\mathit{oh}}
\def\Spacesstack{\mathcal{S}\!\mathit{paces}}
\def\Quotfunctor{\mathrm{Quot}}
\def\Hilbfunctor{\mathrm{Hilb}}
\def\Curvesstack{\mathcal{C}\!\mathit{urves}}
\def\Polarizedstack{\mathcal{P}\!\mathit{olarized}}
\def\Complexesstack{\mathcal{C}\!\mathit{omplexes}}
% \Pic is the operator that assigns to X its picard group, usage \Pic(X)
% \Picardstack_{X/B} denotes the Picard stack of X over B
% \Picardfunctor_{X/B} denotes the Picard functor of X over B
\def\Pic{\mathop{\mathrm{Pic}}\nolimits}
\def\Picardstack{\mathcal{P}\!\mathit{ic}}
\def\Picardfunctor{\mathrm{Pic}}
\def\Deformationcategory{\mathcal{D}\!\mathit{ef}}

\setlength{\emergencystretch}{3em}
\begin{document}
\title[Stacks Project, Tag 032A]{030 --- Cohen structure theorem}
\maketitle
\noindent Copyright (C) 2005--2025 Johan de Jong. Stacks Project authors. Source: \href{https://stacks.math.columbia.edu/tag/032A}{Tag 032A}.

\noindent Editorial difficulty note (not author text). Difficulty H2: The proof constructs coefficient fields or Cohen coefficient rings by successive lifting. It treats mixed characteristic separately and produces the required surjection from a power-series ring, rather than presupposing a coefficient ring..

\noindent Editorial provenance note (not author text). History: excerpt from revision \texttt{a0444\allowbreak 6e57e\allowbreak c1fbc\allowbreak 252a8\allowbreak 71afc\allowbreak ec775\allowbreak 2fb28\allowbreak 07b14}, algebra.tex, lines 44928--45020. Reference presentation was adapted; original-author context and core support blocks were retained or added. The mathematical wording is unchanged. Licensed under GNU FDL 1.2 or later; no invariant sections or cover texts. See ../licenses/COPYING. Context blocks reproduce author definitions and supporting results; they are not additional counted problems.

\begin{definition}
\label{definition-coefficient-ring}
Let $(R, \mathfrak m)$ be a complete local ring.
A subring $\Lambda \subset R$ is
called a {\it coefficient ring} if the following conditions hold:
\begin{enumerate}
\item $\Lambda$ is a complete local ring with maximal ideal
$\Lambda \cap \mathfrak m$,
\item the residue field of $\Lambda$ maps isomorphically to the
residue field of $R$, and
\item $\Lambda \cap \mathfrak m = p\Lambda$, where $p$ is the characteristic
of the residue field of $R$.
\end{enumerate}
\end{definition}

\begin{definition}
\label{definition-cohen-ring}
A {\it Cohen ring} is a complete discrete valuation ring with
uniformizer $p$ a prime number.
\end{definition}

\begin{lemma}
\label{lemma-cohen-rings-exist}
Let $p$ be a prime number.
Let $k$ be a field of characteristic $p$.
There exists a Cohen ring $\Lambda$ with $\Lambda/p\Lambda \cong k$.
\end{lemma}

\begin{proof}
First note that the $p$-adic integers $\mathbf{Z}_p$ form a Cohen ring
for $\mathbf{F}_p$. Let $k$ be an arbitrary field of characteristic $p$.
Let $\mathbf{Z}_p \to R$ be a flat local ring map such that
$\mathfrak m_R = pR$ and $R/pR = k$, see
Lemma \ref{lemma-flat-local-given-residue-field}.
By Lemma \href{https://stacks.math.columbia.edu/tag/05GH}{lemma completion Noetherian (Tag 05GH)} the completion
$\Lambda = R^\wedge$ is Noetherian. Since $\Lambda/p\Lambda = R/pR$
is a field we conclude $\Lambda$ is a
complete Noetherian local ring with maximal ideal $(p)$
(equality and completeness by Lemma \href{https://stacks.math.columbia.edu/tag/05GG}{lemma hathat finitely generated (Tag 05GG)}).
For every $n$ the map $\mathbf{Z}/p^n\mathbf{Z} \to \Lambda/p^n\Lambda = R/p^nR$
is flat (equality by Lemma \href{https://stacks.math.columbia.edu/tag/05GG}{lemma hathat finitely generated (Tag 05GG)} again).
Hence $\mathbf{Z}_p \to \Lambda$ is flat for example by
Lemma \href{https://stacks.math.columbia.edu/tag/0523}{lemma flat module powers (Tag 0523)}. Hence $p$ is a
nonzerodivisor in $\Lambda$. Hence $\Lambda$ has dimension $\geq 1$
(Lemma \href{https://stacks.math.columbia.edu/tag/00KW}{lemma one equation (Tag 00KW)})
and we conclude that $\Lambda$ is regular of dimension $1$, i.e.,
a discrete valuation ring by Lemma \href{https://stacks.math.columbia.edu/tag/00PD}{lemma characterize dvr (Tag 00PD)}.
We conclude $\Lambda$ is a Cohen ring for $k$.
\end{proof}

\begin{lemma}
\label{lemma-cohen-ring-formally-smooth}
Let $p > 0$ be a prime.
Let $\Lambda$ be a Cohen ring with residue field of characteristic $p$.
For every $n \geq 1$ the ring map
$$
\mathbf{Z}/p^n\mathbf{Z} \to \Lambda/p^n\Lambda
$$
is formally smooth.
\end{lemma}

\begin{proof}
If $n = 1$, this follows from
Proposition \href{https://stacks.math.columbia.edu/tag/0322}{proposition characterize separable field extensions (Tag 0322)}.
For general $n$ we argue by induction on $n$.
Namely, if $\mathbf{Z}/p^n\mathbf{Z} \to \Lambda/p^n\Lambda$ is
formally smooth, then we can apply Lemma \href{https://stacks.math.columbia.edu/tag/031L}{lemma lift formal smoothness (Tag 031L)}
to the ring map
$\mathbf{Z}/p^{n + 1}\mathbf{Z} \to \Lambda/p^{n + 1}\Lambda$
and the ideal $I = (p^n) \subset \mathbf{Z}/p^{n + 1}\mathbf{Z}$.
\end{proof}

\begin{lemma}
\label{lemma-flat-local-given-residue-field}
Let $(R, \mathfrak m, k)$ be a local ring. Let $K/k$ be a field
extension. There exists a local ring $(R', \mathfrak m', k')$, a flat local
ring map $R \to R'$ such that $\mathfrak m' = \mathfrak mR'$ and such that
$k'$ is isomorphic to $K$ as an extension of $k$.
\end{lemma}

\begin{proof}
Suppose that $k' = k(\alpha)$ is a monogenic extension of $k$.
Then $k'$ is the residue field of a flat local extension $R \subset R'$
as in the lemma. Namely, if $\alpha$ is transcendental over $k$, then we let
$R'$ be the localization of $R[x]$ at the prime $\mathfrak mR[x]$.
If $\alpha$ is algebraic with minimal polynomial
$T^d + \sum \overline{\lambda}_iT^{d - i}$, then we let
$R' = R[T]/(T^d + \sum \lambda_i T^{d - i})$.

\medskip\noindent
Consider the collection of triples $(k', R \to R', \phi)$, where
$k \subset k' \subset K$ is a subfield,
$R \to R'$ is a local ring map as in the lemma, and
$\phi : R' \to k'$ induces an isomorphism $R'/\mathfrak mR' \cong k'$
of $k$-extensions. These form a ``big'' category $\mathcal{C}$ with morphisms
$(k_1, R_1, \phi_1) \to (k_2, R_2, \phi_2)$
given by ring maps $\psi : R_1 \to R_2$ such that
$$
\xymatrix{
R_1 \ar[d]_\psi \ar[r]_{\phi_1} & k_1 \ar[r] & K \ar@{=}[d] \\
R_2 \ar[r]^{\phi_2} & k_2 \ar[r] & K
}
$$
commutes. This implies that $k_1 \subset k_2$.

\medskip\noindent
Suppose that $I$ is a directed set, and
$((R_i, k_i, \phi_i), \psi_{ii'})$ is a system over $I$, see
Categories, Section \href{https://stacks.math.columbia.edu/tag/002Z}{section posets limits (Tag 002Z)}.
In this case we can consider
$$
R' = \colim_{i \in I} R_i
$$
This is a local ring with maximal ideal $\mathfrak mR'$, and
residue field $k' = \bigcup_{i \in I} k_i$. Moreover, the ring
map $R \to R'$ is flat as it is a colimit of flat maps (and tensor
products commute with directed colimits).
Hence we see that $(R', k', \phi')$ is an ``upper bound'' for the system.

\medskip\noindent
An almost trivial application of Zorn's Lemma would finish the proof
if $\mathcal{C}$ was a set, but it isn't.
(Actually, you can make this work by finding a reasonable bound on the
cardinals of the local rings occurring.)
To get around this problem we choose a well ordering on $K$.
For $x \in K$ we let $K(x)$ be the subfield of $K$ generated
by all elements of $K$ which are $\leq x$.
By transfinite recursion on $x \in K$ we will produce ring maps
$R \subset R(x)$ as in the lemma with residue field extension
$K(x)/k$. Moreover, by construction we will have that
$R(x)$ will contain $R(y)$ for all $y \leq x$.
Namely, if $x$ has a predecessor $x'$, then $K(x) = K(x')[x]$
and hence we can let $R(x') \subset R(x)$ be the local ring extension
constructed in the first paragraph of the proof. If $x$ does not
have a predecessor, then we first set
$R'(x) = \colim_{x' < x} R(x')$ as in the third paragraph
of the proof. The residue field of $R'(x)$ is $K'(x) = \bigcup_{x' < x} K(x')$.
Since $K(x) = K'(x)[x]$ we see that we can use the construction of the
first paragraph of the proof to produce $R'(x) \subset R(x)$.
This finishes the proof of the lemma.
\end{proof}

\begin{theorem}[Cohen structure theorem]
\label{theorem-cohen-structure-theorem}
Let $(R, \mathfrak m)$ be a complete local ring.
\begin{enumerate}
\item $R$ has a coefficient ring (see
Definition \ref{definition-coefficient-ring}),
\item if $\mathfrak m$ is a finitely generated ideal, then
$R$ is isomorphic to a quotient
$$
\Lambda[[x_1, \ldots, x_n]]/I
$$
where $\Lambda$ is either a field or a Cohen ring.
\end{enumerate}
\end{theorem}

\begin{proof}
Let us prove a coefficient ring exists.
First we prove this in case the characteristic of the residue field $\kappa$
is zero. Namely, in this case we will prove by induction
on $n > 0$ that there exists a section
$$
\varphi_n : \kappa \longrightarrow R/\mathfrak m^n
$$
to the canonical map $R/\mathfrak m^n \to \kappa = R/\mathfrak m$.
This is trivial for $n = 1$. If $n > 1$, let $\varphi_{n - 1}$ be given.
The field extension $\kappa/\mathbf{Q}$ is formally smooth by
Proposition \href{https://stacks.math.columbia.edu/tag/0322}{proposition characterize separable field extensions (Tag 0322)}.
Hence we can find the dotted arrow
in the following diagram
$$
\xymatrix{
R/\mathfrak m^{n - 1} &
R/\mathfrak m^n \ar[l] \\
\kappa \ar[u]^{\varphi_{n - 1}} \ar@{..>}[ru] & \mathbf{Q} \ar[l] \ar[u]
}
$$
This proves the induction step. Putting these maps together
$$
\lim_n\ \varphi_n : \kappa \longrightarrow
R = \lim_n\ R/\mathfrak m^n
$$
gives a map whose image is the desired coefficient ring.

\medskip\noindent
Next, we prove the existence of a coefficient ring in the case
where the characteristic of the residue field $\kappa$ is $p > 0$.
Namely, choose a Cohen ring $\Lambda$ with $\kappa = \Lambda/p\Lambda$,
see Lemma \ref{lemma-cohen-rings-exist}. In this case we will prove by
induction on $n > 0$ that there exists a map
$$
\varphi_n :
\Lambda/p^n\Lambda
\longrightarrow
R/\mathfrak m^n
$$
whose composition with the reduction map $R/\mathfrak m^n \to \kappa$
produces the given isomorphism $\Lambda/p\Lambda = \kappa$. This is trivial
for $n = 1$. If $n > 1$, let $\varphi_{n - 1}$ be given.
The ring map $\mathbf{Z}/p^n\mathbf{Z} \to \Lambda/p^n\Lambda$
is formally smooth by Lemma \ref{lemma-cohen-ring-formally-smooth}.
Hence we can find the dotted arrow
in the following diagram
$$
\xymatrix{
R/\mathfrak m^{n - 1} &
R/\mathfrak m^n \ar[l] \\
\Lambda/p^n\Lambda \ar[u]^{\varphi_{n - 1}} \ar@{..>}[ru] &
\mathbf{Z}/p^n\mathbf{Z} \ar[l] \ar[u]
}
$$
This proves the induction step. Putting these maps together
$$
\lim_n\ \varphi_n :
\Lambda = \lim_n\ \Lambda/p^n\Lambda
\longrightarrow
R = \lim_n\ R/\mathfrak m^n
$$
gives a map whose image is the desired coefficient ring.

\medskip\noindent
The final statement of the theorem follows readily. Namely, if
$y_1, \ldots, y_n$ are generators of the ideal $\mathfrak m$,
then we can use the map $\Lambda \to R$ just constructed
to get a map
$$
\Lambda[[x_1, \ldots, x_n]] \longrightarrow R,
\quad x_i \longmapsto y_i.
$$
Since both sides are $(x_1, \ldots, x_n)$-adically complete
this map is surjective by Lemma \href{https://stacks.math.columbia.edu/tag/0315}{lemma completion generalities (Tag 0315)}
as it is surjective modulo $(x_1, \ldots, x_n)$ by
construction.
\end{proof}
\end{document}
